%% kapa: exam=1 round=30 year=2019 period=1 form=A subject=real_estate subject_ko=부동산학원론 seq=22 number=102
%% source: Q-Net 감정평가사 기출문제 (cst00309) artlSeq=5205205 file=extracted/1008/5205205/2019년 제30회 감정평가사 1차시험 문제지/2019년 제30회 감정평가사 1차시험 문제지/2019년도 제30회 감정평가사 1차 1교시 A형.hwp
%% conversion: kordoc 4.24.51 -> tools/gen1.py (auto); review=proofread-sonnet; flags=none
\begin{question}{102}
최유효이용에 관한 설명으로 옳지 않은 것은?
\begin{choices}
  \item 토지이용흡수율 분석은 경제적 타당성 여부판단에 활용되지 않는다.
  \item 인근지역의 용도와는 전혀 다른 데도 불구하고 최유효이용이 되는 경우가 있다.
  \item 중도적 이용에 할당되고 있는 부동산을 평가할 때는 토지와 개량물을 같은 용도로 평가해야 한다.
  \item 단순히 최고의 수익을 창출하는 잠재적 용도가 아니라 적어도 그 용도에 대한 유사부동산의 시장수익률과 동등 이상의 수준이 되어야 한다.
  \item 투기적 목적으로 사용되고 있는 토지에 대한 최유효이용분석에 있어서는 특정한 용도를 미리 상정해서는 안되며 미래 사용에 대한 일반적 유형을 상정해야 한다.
\end{choices}
\end{question}
\answer{1}
